\documentclass[10pt]{article}

\usepackage[utf8]{inputenc}

\usepackage[T1]{fontenc}

\usepackage{amsmath}

\usepackage{amsfonts}

\usepackage{amssymb}

\usepackage[version=4]{mhchem}

\usepackage{stmaryrd}

\usepackage{bbold}

\usepackage{mathrsfs}



\begin{document}

тицы. Б. з. Д. является основным понятием квантовой теории кристаллич. твердых тел.



Математически корректно функция $E(\boldsymbol{k})$ определяется следующим образом. Оператор $T_{\boldsymbol{a}}$ сдвига на вектор $\boldsymbol{a}$ решетки $\Lambda$ :



\[

\left(T_{a} \psi\right)(x)=\psi(x-a)

\]



имеет в гильбертовом пространстве $H=L_{2}\left(\mathbb{R}^{N}\right)$ спектральное разложение



\[

T_{a}=\int e^{i k a} d E_{k}

\]



где $\boldsymbol{k}$ - квазиимпульс, характеризующий неприводимое представление группы трансляций. Оператор Шредингера квантовой частицы в периодич. потенциале $-\Delta+u(\boldsymbol{x})$, $u(\boldsymbol{x}+\boldsymbol{a})=u(\boldsymbol{x}), \boldsymbol{a} \in \Lambda$, коммутирует со всеми операторами $T_{a}, \boldsymbol{a} \in \Lambda$, и поэтому в гильбертовом пространстве $H$ имеет спектральное разложение



\[

-\Delta+u(x)=\int \varepsilon(\boldsymbol{k}) d E_{\boldsymbol{k}}

\]



Определенная таким образом функция $\varepsilon(\boldsymbol{k})$ и называется законом дисперсии Блоха.



Функция $\varepsilon(\boldsymbol{k})$ обладает рядом симметрий. Прежде всего, для самосопряженного оператора Шредингера эта функция четная: $\varepsilon(-\boldsymbol{k})=\varepsilon(\boldsymbol{k})$. Поскольку квазиимпульс $\boldsymbol{k}$ определен с точностью до произвольного вектора $\boldsymbol{b}$ обратной решетки $\Lambda^{*}, \varepsilon(\boldsymbol{k})$ является периодич. функцией с периодами $\boldsymbol{b}$, $\boldsymbol{\varepsilon}(\boldsymbol{k}+\boldsymbol{b})=\boldsymbol{\varepsilon}(\boldsymbol{k}), \boldsymbol{b} \in \Lambda^{*}$, и поэтому может рассматриваться как функция на фундаментальной области обратной решетки, называемой Бриллюэна зоной. Если потенциал $u(\boldsymbol{x})$ инвариантен относительно пространственной группы $G$, подгруппой к-рой являетсся решетка сдвигов $\Lambda$, то $\varepsilon(r \boldsymbol{k})=\varepsilon(\boldsymbol{k}), r \in R$, где $R$ - точечная группа $G$, то есть закон дисперсии инвариантен относительно точечной группы.



$\varepsilon(\boldsymbol{k})$ является многозначной аналитич. функщией компонент квазиимпульса. Различные ветви этой функции $\varepsilon_{n}(\boldsymbol{k})$, $n \in \mathbb{N}$, называются эне ргетическими зонами. Расположение точек ветвления в ряде случаев может быть определено из соображений симметрии. Напр., точка $\boldsymbol{k}$ является точкой ветвления $p$-го порядка, если группа вектора $\boldsymbol{k}$, к-рая есть $R_{\boldsymbol{k}}=\{r \mid r \boldsymbol{k}=k\}$, имеет $p$-мерное представление. Изоэнергетич. поверхность $\Gamma_{\boldsymbol{a}}=\{\boldsymbol{k} \mid \varepsilon(\boldsymbol{k})=\boldsymbol{a}-$ const $\}$ не может содержать кусков поверхностей, параметризованных рациональными функциями, в частности куски плоскостей, прямых цилиндров, конусов, параболоидов, гиперболоидов. При $N>1$ области значений функций $\varepsilon_{n}(\boldsymbol{k})$ всегда перекрываются при достаточно больших $\varepsilon$. При $N=1$ последовательность функций $\varepsilon_{n}(k), k \in[0, \pi / a], n \in \mathbb{N}$, состоит из поочередно монотонно возрастающих и убывающих функций, области значений этих функций могут лишь попарно касаться друг друга, и если они касаются всегда, за исключением конечного множества зон, то соответствующий оператор Шредингера называется кон е ч о о о н ны м.



Для гладких потенциалов спектр оператора Шредингера является абсолютно непрерывным. Плотность состояний спектра $p(E)$ выражается через закон дисперсии $\varepsilon(\boldsymbol{k})$ по формуле



\[

\begin{aligned}

& p(E)=(2 \pi)^{-N} \int \delta(\varepsilon(\boldsymbol{k})-E) d \boldsymbol{k}= \\

& =(2 \pi)^{-N} \int_{\varepsilon(\boldsymbol{k})=E}\left|\operatorname{grad}_{\boldsymbol{k}} \varepsilon(\boldsymbol{k})\right|^{-1} d s .

\end{aligned}

\]



Особенности плотности состояний определяются критич. точками закона дисперсии, где $\operatorname{grad}_{\boldsymbol{k}} \varepsilon(\boldsymbol{k})=0$. Характер особенности $p(E)$ в невырожденной критич. точке определяется индексом критич. точки (то есть числом отрицательных собственных значений матрицы $\left\|\partial^{2} \varepsilon / \partial k_{i} \partial k_{j}\right\|$ в критич. точке) и размерностью пространства $N$. Поскольку функцию $\varepsilon(\boldsymbol{k})$ можно считать заданной на зоне Бриллюэна или, что то же самое, на $N$-мерном торе, числа $\mu_{\lambda}$ критич. точек индекса $\lambda$ этой функции удовлетворяют н е р а в е н с т в а м М о р с а $\mu_{\lambda} \geqslant b_{\lambda}$, где $b_{\lambda}$ - числа Бетти тора. В частности, для $N=3$ плотность состояний спектра $p(E)$ может иметь только корневые особенности, а $b_{0}=1, b_{1}=3, b_{2}=3, b_{3}=1$.



С помошью Б. з. д. можно рассчитать асимптотич. методами спектр оператора Шредингера, учитывающего помимо периодического также произвольное слабое поле (слабое в том смысле, что характеристики поля мало меняются на длине векторов элементарных трансляций). Напр., для слабого однородного магнитного поля $\boldsymbol{H}$ спектр оператора Шредингера $\varepsilon\left(n, K_{H}\right)$, где $n \in \mathbb{N}, K_{H}-$ проекция квазиимпульса на направление магнитного поля $\boldsymbol{v}=\boldsymbol{H} /|\boldsymbol{H}|$, определяется неявным образом при помощи уравнения



\[

S\left(\varepsilon, K_{H}\right)=(e \boldsymbol{H} / \hbar c) n

\]



где $\hbar-$ постоянная Планка, $c$ - скорость света, $e-$ заряд электрона, а $S\left(\varepsilon, K_{H}\right)-$ площадь сечения изоэнергетич. поверхности плоскостью $(\boldsymbol{k}, \boldsymbol{v})=\boldsymbol{K}_{\boldsymbol{H}}$.



Название термина связано с именем Ф. Блоха (F. Bloch), К-рый в 20-е гг. сформулировал основы Квантовой теории кристаллич. твердых тел.



Лит.: [1] Р и д М., С а й м о н Б., Методы современной математической физики, пер. с англ., т. 4, Анализ операторов, М., 1982. Е.Д. Белоколос.



БЛОХА ТЕОРЕМА - фундаментальная теорема квантовой теории твердого тела, устанавливающая вид волновой функции электрона, находящегося в поле с периодическим потенциалом $U$, в частности в кристаллической решетке. Сформулирована Ф. Блохом (F. Bloch) в 1929. Б. т. утверждает, что если потенциал $U(\boldsymbol{r})(\boldsymbol{r}-$ пространственная координата) - функция с периодом $\boldsymbol{a}$ кристаллич. решетки: $U(\boldsymbol{r}+\boldsymbol{a})=U(\boldsymbol{r})$, где $\boldsymbol{a}=n_{1} \boldsymbol{a}_{1}+n_{2} \boldsymbol{a}_{2}+n_{3} \boldsymbol{a}_{3} ; \boldsymbol{a}_{1}, \boldsymbol{a}_{2}, \boldsymbol{a}_{3}$ - основные (базисные) векторы решетки; $n_{1}, n_{2}, n_{3}$ - целые числа, то решения $\psi(\boldsymbol{r})$ одноэлектронного Шредингера уравнения (адиабатическое приближение)



\[

\hat{H} \psi(\boldsymbol{r}) \equiv\left[-\frac{\hbar^{2}}{2 m} \nabla^{2}+U(\boldsymbol{r})\right] \psi(\boldsymbol{r})=\mathscr{E} \psi(\boldsymbol{r})

\]



( $\mathscr{E}$ - энергия частицы) имеют вид:



\[

\psi_{\boldsymbol{k}}(\boldsymbol{r})=e^{i k r_{k}(\boldsymbol{r}) .}

\]



Здесьь $\boldsymbol{k}$ - волновой вектор, характеризующий состояние электрона, $u_{k}-$ периодич. функция с периодом решетки, $m$ - масса электрона. Б. т. является следствием трансляционной инвариантности кристаллич. решетки. Если $\psi(\boldsymbol{r})-$ решение уравнения (1), соответствующее стационарному состоянию электрона с энергией $\mathscr{E}$, то $\psi(\boldsymbol{r}+\boldsymbol{a})$ также является его решением, причем $\psi(\boldsymbol{r}+\boldsymbol{a})=C_{a} \psi(\boldsymbol{r})$, где $C_{a}= \pm 1$



Если стационарному состоянию с энергией $\mathscr{E}$ соответствует несколько различных волновых функций $\psi(\boldsymbol{r})$ (то есть состояние с энергией $\mathscr{E}$ - вырожденное), то волновая функция $\psi(\boldsymbol{r}+\boldsymbol{a})$ является линейной комбинацией всех собственных функций $\psi(\boldsymbol{r})$, отвечающих вырожденному уровню ๕. В этом случае $C_{\boldsymbol{a}}=e^{i \boldsymbol{k} a}$. Таким образом, в случае вырождения имеем:



\[

\psi(\boldsymbol{r}+\boldsymbol{a})=e^{i k a} \psi(\boldsymbol{r}) .

\]



Функции, удовлетворяюшие условию (3) (у с ло в и ю Блоха), называются функциями Блоха.



Подставляя (2) в уравнение Шредингера (1), получим уравнение для $u_{k}(\boldsymbol{r})$ :



\[

\left[-\frac{\hbar^{2}}{2 m}(i \nabla+\boldsymbol{k})^{2}+U(\boldsymbol{r})\right] u_{\boldsymbol{k}}(\boldsymbol{r})=\mathscr{E}_{\boldsymbol{k}} U(\boldsymbol{r}),

\]



которое имеет бесконечный ряд решений $u_{s k}(\boldsymbol{r})$. Индекс $s$ нумерует решения при заданном $\boldsymbol{k}$. Волновым функциям (2) при заданном $\boldsymbol{k}$, таким образом, соответствуют дискретные значения энергии: $\varepsilon=\mathscr{E}_{s}(\boldsymbol{k}), s=1,2, \ldots$ Индекс $s-$ номер энергетич. зоны; зависимость $\mathscr{E}_{S}$ от $\boldsymbol{k}$ при фиксированном $s$ называется Блоха законом дисперсии частицы в $s$-й зоне.



Лит.: [1] 3 а й м а н Дж., Принципы теории твердого тела, пер. с англ., М., 1974; [2] Л и фш иц Е. М., П и тае вски й Л. П., Статистическая физика, ч. 2, М., 1978.



В. М. Винокур





\end{document}